\input{C:/Users/pif-m/OneDrive/Desktop/AI/Kurser/DYN 1/DYN_Eksamen_Svar/preamble.tex}

\renewcommand{\hdtitel}{MEK2 Læseguide --- Lektion 1: Spænding og tøjning}

\title{Lektion 1: Spænding og tøjning (GG 1.4, 1.7, 1.8)}
\author{MEK2}
\date{}

\begin{document}

\maketitle

% === LÆRINGSMÅL ===
\begin{spgbox}
\textbf{Læringsmål:} Efter denne lektion kan du
\begin{itemize}
  \item Beregne normalspænding $\sigma$ og normal tøjning $\varepsilon$ fra kraft, areal og deformation
  \item Identificere centroiden og forstå betydningen for uniform lastfordeling
  \item Anvende Hookes lov $\sigma = E\varepsilon$ inden for det lineære regime
  \item Beregne lateral tøjning via Poissons forhold $\nu$
  \item Beregne forskydningsspænding $\tau$ i bolte og samlinger
  \item Skelne mellem single shear og double shear
  \item Beregne lejespænding (bearing stress) $\sigma_b$ fra projiceret areal
  \item Anvende Hookes lov i forskydning: $\tau = G\gamma$ med $G = \frac{E}{2(1+\nu)}$
\end{itemize}
\end{spgbox}

% === REFERENCER ===
\begin{refbox}
\textbf{Pensum:} Goodno \& Gere: Mechanics of Materials, SI ed., enhanced 9th ed., kapitel 1.4, 1.7, 1.8
\end{refbox}

\clearpage

% === KERNEBEGREBER ===
\begin{nembox}

\textbf{Ni kernebegreber fra slidesne:}

\begin{enumerate}
  \item \textbf{Normalspænding} $\sigma$ = kraft pr. areal: $\sigma = \frac{P}{A}$ [Pa]. Måler trykbelastning vinkelret på tværsnit.

  \item \textbf{Normal tøjning} $\varepsilon$ = relativ forlængelse: $\varepsilon = \frac{\delta}{L}$ [dimensionsløs]. Hvor $\delta$ er den absolutte længdeændring.

  \item \textbf{Youngs modul} $E$ [Pa] = stivhedskonstant i normalretning. Beskriver modstand mod elastisk deformation.

  \item \textbf{Poissons forhold} $\nu$ [dimensionsløs] = relation mellem lateral og aksial tøjning: $\nu = -\frac{\varepsilon_{\text{lat}}}{\varepsilon_{\text{aks}}}$. Typisk 0,2--0,35.

  \item \textbf{Centroide} $(\bar{x}, \bar{y})$ = geometrisk tyngdepunkt. Resultanten af uniform normalspænding ligger her: $\bar{x} = \frac{\int x \, dA}{A}$, $\bar{y} = \frac{\int y \, dA}{A}$.

  \item \textbf{Forskydningsspænding} $\tau$ = kraft pr. areal parallelt med overflade: $\tau_{\text{aver}} = \frac{V}{A}$ [Pa]. Karakteristisk for bolte og samlinger.

  \item \textbf{Lejespænding} $\sigma_b$ = normalspænding fra lateral kontakt mellem bolt og hul: $\sigma_b = \frac{F_b}{A_b}$ [Pa]. Udregnes med projiceret areal $A_b = t \cdot d_b$.

  \item \textbf{Forskydningstøjning} $\gamma$ [rad] = ændring i vinkel mellem vinkelrette planer. Positivt når vinklen mellem ``positive'' flader mindskes.

  \item \textbf{Forskydningsmodul} $G$ [Pa] = stivhedskonstant i forskydning: $G = \frac{E}{2(1+\nu)}$. Relaterer $\tau$ og $\gamma$ via Hookes lov: $\tau = G\gamma$.

\end{enumerate}

\end{nembox}

\clearpage

% === FORMLER ===
\begin{eksambox}

\textbf{Alle ni ligninger fra slidesne (GG 1.4, 1.7, 1.8):}

\vspace{8pt}

\textbf{Normalspænding og -tøjning:}
\begin{align}
  \sigma &= \frac{P}{A} \quad \text{(1--6)} \quad [\text{Pa}] \quad \text{hvor } P \text{ = kraft [N], } A \text{ = tværsnitsareal [m}^2\text{]} \\
  \varepsilon &= \frac{\delta}{L} \quad \text{(1--7)} \quad [\text{-}] \quad \text{hvor } \delta \text{ = forlængelse [m], } L \text{ = oprindelig længde [m]} \\
  \bar{x} &= \frac{\int x \, dA}{A}, \quad \bar{y} = \frac{\int y \, dA}{A} \quad \text{(1--9a,b)} \quad \text{centroide-formler}
\end{align}

\textbf{Hookes lov og Poissons forhold:}
\begin{align}
  \sigma &= E\varepsilon \quad \text{(1--12)} \quad \text{hvor } E \text{ = Youngs modul [Pa]} \\
  \nu &= -\frac{\varepsilon_{\text{lat}}}{\varepsilon_{\text{aks}}} \quad \text{(1--13)} \quad \text{hvor } -1 \le \nu \le 0.5 \text{ (typisk } \sim 0.2\text{--}0.35\text{)}
\end{align}

\textbf{Forskydningsspænding og -tøjning:}
\begin{align}
  \sigma_b &= \frac{F_b}{A_b} \quad \text{(1--15)} \quad A_b = t \cdot d_b \quad \text{(lejespænding, projiceret areal)} \\
  \tau_{\text{aver}} &= \frac{V}{A} \quad \text{(1--16)} \quad \text{hvor } V = \frac{P}{2} \text{ (double shear) eller } V = P \text{ (single shear)} \\
  \tau_1 &= \tau_2 \quad \text{(1--17)} \quad \text{(tangentialspændinger på modstående flader er ens)} \\
  \tau &= G\gamma \quad \text{(1--18)} \quad \text{hvor } G \text{ = forskydningsmodul [Pa], } \gamma \text{ = forskydningstøjning [rad]} \\
  G &= \frac{E}{2(1+\nu)} \quad \text{(1--19)} \quad \text{relation mellem } E \text{, } G \text{ og } \nu
\end{align}

\textbf{Vigtige noter:}
\begin{itemize}
  \item Hookes lov ($\sigma = E\varepsilon$ og $\tau = G\gamma$) gælder kun i det \emph{lineære regime} før flydning.
  \item Aksial last gennem centroide giver rent træk/tryk uden bøjningsmoment.
  \item Lejespænding kan blive kritisk hvis tykkelsen $t$ er lille.
\end{itemize}

\end{eksambox}

\clearpage

% === HUSKEORD ===
\begin{opfbox}

\textbf{Huskeord og mnemoteknik:}

\begin{itemize}
  \item \textbf{$\sigma$ og $\varepsilon$:} ``Spænding = kraft divideret areal, tøjning = længdeændring divideret oprindelig længde.''

  \item \textbf{Centroiden:} ``Kraftresultanten ligger i centroiden --- det geometriske tyngdepunkt.''

  \item \textbf{Hookes lov lineær:} ``$\sigma = E\varepsilon$ er lineær og gælder kun før flydning. Kontroller diagrammet!''

  \item \textbf{Poissons forhold negativ:} ``$\nu = -\frac{\varepsilon_{\text{lat}}}{\varepsilon_{\text{aks}}}$ er negativ fordi når materiale strekkes aksielt (positivt $\varepsilon_{\text{aks}}$), bliver det smalere (negativt $\varepsilon_{\text{lat}}$).''

  \item \textbf{Double vs. single shear:} ``En bolt i en samling kan være i double shear (belastet fra begge sider, $V = P/2$) eller single shear (belastet fra én side, $V = P$).''

  \item \textbf{Lejespænding:} ``Brug \emph{projiceret areal} $A_b = t \cdot d_b$ (tykkelse gange bolt-diameter), ikke cirkelareal!''

  \item \textbf{Forskydningsmodul:} ``$G = \frac{E}{2(1+\nu)}$ --- G kan beregnes fra E og Poissons forhold. For stål typisk $G \approx 80$ GPa.''

  \item \textbf{Forskydningstøjning $\gamma$:} ``Måles i radianer. $\gamma > 0$ når vinklen mellem positive flader mindskes. Små vinkler: $d = t \tan(\gamma) \approx t \cdot \gamma$.''

\end{itemize}

\end{opfbox}

\end{document}
